% Generated draft: requires mathematical and source review.
\documentclass[11pt]{article}
\usepackage[margin=25mm]{geometry}
\usepackage{fontspec}
\setmainfont{Latin Modern Roman}
\usepackage{amsmath,amssymb,mathtools,unicode-math}
\setmathfont{Latin Modern Math}
\usepackage{enumitem,needspace,xurl,hyperref}
\setlength{\parindent}{0pt}
\setlength{\parskip}{5pt}
\setlength{\emergencystretch}{3em}
\newcommand{\N}{\mathbb{N}}
\newcommand{\Z}{\mathbb{Z}}
\newcommand{\Q}{\mathbb{Q}}
\newcommand{\R}{\mathbb{R}}
\newcommand{\C}{\mathbb{C}}
\newcommand{\F}{\mathbb{F}}
\newcommand{\A}{\mathbb{A}}
\newcommand{\PP}{\mathbb{P}}
\newcommand{\E}{\mathbb{E}}
\newcommand{\eps}{\varepsilon}
\newcommand{\dd}{\,\mathrm{d}}
\begin{document}
\section*{UnsolvedMath ID: 801076}
\textbf{Source problem code:} catalogue entry
\subsection{Definitions and mathematical statement}
The supplied record does not define the family or the set \(J^+_{a,c}\), or specify the parameter domain. The following formulation is therefore conditional on the complex quadratic H\'enon convention
\[
H_{a,c}(x,y)=(x^2+c-ay,x),\qquad (x,y)\in\C^2,\quad (a,c)\in\C^*\times\C,
\]
where \(\C^*=\C\setminus\{0\}\). These conventions require verification against the source. Under these conventions, define
\[
K^+_{a,c}=\left\{z\in\C^2:\sup_{n\geq0}\|H_{a,c}^{\,n}(z)\|<\infty\right\},\qquad
J^+_{a,c}=\partial K^+_{a,c},
\]
where the boundary is taken in \(\C^2\), and let
\[
d(a,c)=\dim_{\mathrm H}J^+_{a,c}
\]
be the Hausdorff dimension for the Euclidean metric on \(\C^2\cong\R^4\).

The requested problem is to determine an asymptotic expansion of \(d(a,c)\) as \((a,c)\to(0,0)\). The source does not specify the scale of the expansion, its remainder, or restrictions on the relative rates at which the two parameters approach zero.

In particular, the source asks whether there is a lower bound
\[
d(a,c)\geq 3+h(a),\qquad h(a)>0,
\]
uniform for \(c\) in a neighborhood of zero. A formulation with one common neighborhood would be: do there exist \(\delta,\varepsilon>0\) and a function
\[
h:\{a\in\C:0<|a|<\delta\}\longrightarrow(0,\infty)
\]
such that
\[
\forall a\ (0<|a|<\delta)\quad\forall c\ (|c|<\varepsilon),\qquad d(a,c)\geq3+h(a)?
\]
A weaker interpretation permits the neighborhood of \(c=0\) to depend on \(a\), replacing this condition by
\[
\forall a\ (0<|a|<\delta)\quad\exists\varepsilon_a>0\quad\forall c\ (|c|<\varepsilon_a),\qquad d(a,c)\geq3+h(a).
\]
The record does not determine which quantifier order is intended. Excluding \(a=0\) in these conditional formulations reflects the invertible H\'enon parameter domain assumed above; the source does not explicitly specify whether the degenerate parameter itself is included.
\subsection{Short English statement}
Determine how the Hausdorff dimension of the forward Julia set \(J^+_{a,c}\) behaves as \((a,c)\to(0,0)\). In particular, ask whether it exceeds \(3\) by a positive quantity depending only on \(a\), uniformly for \(c\) near zero. The family convention, parameter domain, treatment of \(a=0\), and precise uniformity require source verification.
\subsection{Sources}
\begin{itemize}
\item Local UnsolvedMath registry, permanent ID \(801076\), exact source problem code \texttt{catalogue entry}, title: \emph{Henon Maps: Open Problems --- Question 28}. The supplied registry labels the problem open; this is a supplied status, not an independently verified literature assessment.
\item Berger and Bedford, \emph{Henon Maps: A List of Open Problems} (2023), Question 28. Supplied source URL: \url{https://arxiv.org/abs/2312.03907}. The registry attributes the problem to Berger and Bedford, with proposed year 2023.
\item Supplied provenance records access on 2026-10-16 and extraction by \texttt{arxiv\_latex}. This is supplied provenance; no independent source visit or verification is claimed. Supplied upstream revision: \texttt{d371ee917ac9647fbb06cd6fc529deec6f3dda82}.
\end{itemize}
\end{document}
